\section{Discussion}
\label{sec:discussion}

Symmetric bodies walk further than tree bodies with the same number of hinges, but mirroring is not required for this. Tree bodies selected for the same resting shape do as well, and so do tied bodies, which have neither mirroring nor the symmetric shape. Mirroring may help mainly by making low, sprawled bodies more likely, since tree bodies with that shape catch up. It cannot be the only route, since tied bodies get there another way. We do not yet know how: they have more leg hinges, but matching tree bodies on leg hinges does not reproduce their score, even though adjusting for leg hinges by regression shrinks their lead (from $+0.29$ to $+0.10$~m under \bb). The matching experiment is the more direct test, so we do not treat leg count as the explanation.

The lead builds during learning and is absent after the first evaluations, so it is not a head start of bodies that move well with any controller. These bodies gain more from training. The ratio to tree falls slightly after its early peak while the absolute gap keeps growing, so whether the lead shrinks with training depends on the measure.

Symmetric bodies phase-lock, and locked gaits score higher, so it is natural to think that the lock is what makes them good. Three results argue against this for trained gaits: perturbing left--right timing costs no more than perturbing other timing, breaking most of the lock by detuning costs no more than the control, and tied bodies reach symmetric scores without locking. Locking also goes with flatter bodies, so it may reflect the same shape differences that carry the lead. Whether the lock helps during learning remains untested; training symmetric bodies with the two sides held away from lock would test it.

For work that compares body encodings, the practical point is that an encoding changes which bodies are likely. Before crediting a regularity with an advantage, it is worth comparing against bodies that match the shape and duplication the encoding produces, without the regularity itself.

The study has several limitations. The effects are modest and scores are skewed: a randomly chosen symmetric body beats a randomly chosen tree body only 58\% of the time, the \dann\ gap has an interval that includes zero, and tests are uncorrected for multiple comparisons. The evolution experiment has 20 seeds per arm, too few to resolve a symmetric lead of this size given its variance across seeds. Our conclusion is that mirroring is not required, not that it contributes nothing, since the intervals on the comparisons with symmetric bodies allow differences as large as the lead. The six shape measures were chosen because they predicted score in pilot runs, so shape matching partly selects for good bodies by design, and the result shows that these measures are sufficient rather than that they are the cause. Tree bodies also keep more modules than symmetric bodies after matching (26 vs 19). In evolution, tied bodies grow beyond the module limit, although the \bb\ result survives removing those runs. Finally, we study one task on flat ground, with one hinge setting and one set of modules, and the fixed-body results use random bodies rather than evolved ones.
